%\subsection{Limitation on Long-Distance Knowledge Retrieval}
Our InSemRAG involves a locality assumption where we conjecture that the complete referent knowledge is located within the neighborhood of the retrieved chunk.
This influences our SPC design, where we repair the semantically fragmented chunk by backtracking within its neighborhood.
In some cases, referent knowledge may be located beyond the neighborhood of the retrieved chunk.
Thus, simple backtracking to neighboring chunks may fail to recover missing semantics.
Note that the SPC module only intentionally searches the neighboring information during backtracking under the following assumptions:
\begin{itemize}
    \item locality assumption where evidence is only utilized to support local reasoning,
    \item expansion of the backtracking window may unintentionally lead to the retrieval of irrelevant information, and
    \item SPC is utilized as an evidence repair mechanism, instead of document-level co-reference resolution.
\end{itemize}
When backtracking fails to solve the semantic fragmentation, InSemRAG relies on the second level of auditing (coverage audit) to trigger new retrieval round.